% !TeX program = xelatex
% UTF-8 transcription; source wording and doses are retained, not corrected.
% pic01.jpg = 20261012_180426792.jpg (page 1)
% pic02.jpg = 20261012_180433314.jpg (page 2)
% pic03.jpg = 20261012_180439481.jpg (page 3)
% pic04.jpg = 20261012_180444764.jpg (page 4)
% pic05.jpg = 20261012_180449540.jpg (page 5)
% Standalone XeLaTeX document; also usable with \input in a XePersian document.
\ifdefined\settextfont
  \def\EpistaxisStandalone{0}
\else
  \def\EpistaxisStandalone{1}
  \documentclass[12pt,a4paper]{article}
  \usepackage[margin=2cm]{geometry}
  \usepackage{amsmath,amssymb}
  \usepackage{xepersian}
  \IfFontExistsTF{Amiri}{\settextfont{Amiri}}{\settextfont{FreeSerif}}
  \setlatintextfont{FreeSerif}
\fi
\providecommand{\en}[1]{\lr{#1}}
\newcommand{\epipage}[1]{\clearpage\section*{صفحهٔ #1}
  {\footnotesize\en{Subject: \qquad Date: / /}\par}}
\newcommand{\epiprint}{\par{\footnotesize\en{YASHA}\par}}
\newcommand{\epiunclear}[1]{[خوانش نامطمئن: #1]}
\ifnum\EpistaxisStandalone=1 \begin{document}\fi

\epipage{۱}
\section*{\en{Chapter 244 -- Epistaxis}}
تینتالی
\begin{itemize}
\item ۶۰٪ مردم در طول عمر خود دچار اپیستاکسی می‌شوند.
\item کدام بیماران اپیستاکسی به اورژانس می‌آیند؟ زیر ۱۰ سال / بالای ۵۰ سال.
\end{itemize}
\subsection*{علل شایع اپیستاکسی}
\begin{itemize}
\item \en{Digital Trauma}.
\item انحراف سپتوم.
\item خشکی مخاط بینی.
\item رینوسینوزیت.
\item نئوپلاسم.
\item تحریک شیمیایی (کوکائین استنشاقی، آمونیاک).
\end{itemize}
\subsection*{علل سیستمیک که خطر خونریزی را بالا می‌برند}
\begin{itemize}
\item نارسایی مزمن کلیوی.
\item الکلیسم.
\item فشار خون.
\item مالفورماسیون عروقی مثل \en{HHT} (هرِدیتری هموراژیک تلانژکتازی).
\item کوآگولوپاتی‌ها.
\item مصرف آنتی‌کوآگولان.
\item بیماری فون ویلیبراند (\en{vWD}).
\item هموفیلی.
\end{itemize}
\epiprint
{\footnotesize شمارهٔ دست‌نویس: ۱.\par}

\epipage{۲}
\subsection*{آناتومی و پاتوفیزیولوژی}
\begin{itemize}
\item شاخهٔ سوپریور \en{labial} از شریان \en{facial}
  + شاخهٔ انتریور اتموئیدال
  + شاخهٔ ترمینال شریان اسفنوپالاتین
  $\Leftarrow$ شبکهٔ کیسل‌باخ در انتریور نازال سپتوم.
\item محل ۹۰٪ موارد اپیستاکسی است.
\item در رینوسکوپی قدامی قابل مشاهده است.
\item شایع‌ترین محل برای خونریزی‌های خلفی: شریان اسفنوپالاتین؛
  شاخهٔ انتهایی شریان اینترنال ماگزیلاری؛ شاخهٔ اکسترنال کاروتید.
\end{itemize}
\subsection*{یافته‌ها بالینی}
\begin{itemize}
\item در شرح حال: الکل، دارو، \en{NSAID}، مواد مخدر، ضدانعقاد، ضدپلاکت؛
  سابقهٔ تروما به بینی، جراحی قبلی بینی؛ شرح حال کوآگولوپاتی.
\item برای معاینه بیمار باید به پرهیز \en{Sniffing} توصیه شود.
\end{itemize}
\subsection*{تشخیص}
\begin{itemize}
\item اپیستاکسی قدامی از خلفی را افتراق دهید.
\item در اورژانس موارد خلفی را زمانی تشخیص می‌دهیم که دیگر موارد شکست بخورد.
\item شواهد به نفع خونریزی خلفی:
  \begin{itemize}
  \item بیمار مسن با کوآگولوپاتی ارثی یا اکتسابی.
  \item رؤیت خون حجیم در نازوفارنکس.
  \item خونریزی دوطرفه.
  \item عدم پاسخ به پک قدامی.
  \end{itemize}
\end{itemize}
\epiprint
{\footnotesize شمارهٔ دست‌نویس: ۲.\par}

\epipage{۳}
\subsection*{بررسی آزمایشگاهی}
\begin{itemize}
\item نیاز نیست مگر \epiunclear{هیچ بیمار به همودینامیک ناپایدار آن باشد}
  یا خونریزی به سختی کنترل شود.
% The clause after "مگر" is indistinct; its visible reading is retained above.
\item اگر آزمایش خواستیم:
  \begin{itemize}
  \item \en{CBC}.
  \item پنل انعقادی.
  \item \en{BG} و کراس‌مچ.
  \item پنل کامل متابولیک (ارزیابی کبدی).
  \end{itemize}
\end{itemize}
\subsection*{درمان اولیه}
\begin{itemize}
\item \en{Rapid primary survey}: در اختلال در راه هوایی یا همودینامیک بیمار.
\item \en{IV line} در بیمار با خونریزی شدید.
\item کراس‌مچ اگر همودینامیک ناپایدار است.
\item رفع کوآگولوپاتی به آنالیز آزمایش‌ها با دادن فرآوردهٔ خونی.
\item پایین آوردن سریع \en{BP} در اپیستاکسی حاد توصیه نمی‌شود.
\item زمانی \en{BP} را پایین می‌آوریم (کم و آهسته) که نیاز به \en{packing} یا
  مداخلهٔ جراحی پیدا کند؛ کاهش فشار هیدرواستاتیک به تشکیل لخته کمک می‌کند.
\end{itemize}
\subsection*{فشار مستقیم بر روی بینی}
\begin{itemize}
\item در ابتدا می‌خواهیم که فین کند تا مخاط‌های زائد و لختهٔ کوچک‌ها پاک شود.
\item از یک وازوکانستریکتور موضعی استفاده می‌کنیم:
  \begin{itemize}
  \item \en{Oxymetazolin 0.5\%}.
  \item فنیل‌افرین.
  \end{itemize}
\item پرهیز \en{Sniffing}، رو به جلو \en{lean forward}.
\item \en{10--15 min} بر روی بینی فشار وارد کند و با دهان تنفس کند.
\end{itemize}
\epiprint
{\footnotesize شمارهٔ دست‌نویس: ۳.\par
لبهٔ پایین تصویر، از صفحهٔ زیرین: خونریزی دوطرفه؛ عدم پاسخ به پک قدامی؛ ۲.\par}

\epipage{۴}
\subsection*{کاتتریزاسیون \en{chemical}}
\begin{itemize}
\item اگر ۲ نوبت تلاش با فشار دادن موفق نبود.
\item \en{Silver Nitrate} قدم بعدی برای درمان اولیه است.
\item در ابتدا گاز آغشته به نسبت \en{1:1}، \en{Oxy 0.05\%} و لیدو ۴٪
  (۳ گاز) برای بی‌حسی ناحیه.
\item اگر شریان خونریزی‌دهنده رؤیت شده است از این روش استفاده می‌کنیم.
\item وقتی رگ شناسایی شد، چوب سیلور نیترات را به پروگزیمال رگ می‌زنیم.
\item باید در محیط بدون خون باشد، وگرنه جواب نمی‌دهد.
\item چند ثانیه آن را می‌گذاریم بعد برمی‌داریم.
\item هرگز نباید دوطرفه انجام شود.
\item تا ۴–۶ هفته بعد از یک اقدام، نباید دوباره تکرار گردد.
\end{itemize}
\subsection*{ژل‌ها ترومبوژنیک و ترانگزامیک اسید}
\begin{itemize}
\item \en{200 mg} از محلول \en{100 mg/ml}، \en{TXA} اتمایز شود در ناستریل.
\item \en{500 mg TXA} در \en{5 cc NS} اتمایز در ناستریل.
\item تامپون آغشته به \en{500 mg TXA} سپس داخل ناستریل.
% The handwritten quantities above are transcribed verbatim.
\end{itemize}
\subsection*{پک قدامی بینی}
\begin{itemize}
\item اگر فشار مستقیم، کوتر شیمیایی، وازوکانستریکتور جواب نداد و
  ژل و فوم ترومبوژنیک در دسترس نبود.
\end{itemize}
\epiprint
{\footnotesize شمارهٔ دست‌نویس: ۴.\par}

\epipage{۵}
\subsection*{تعیین و فالوآپ}
\begin{itemize}
\item خونریزی کنترل، همودینامیک پایدار: بیش از ۱ ساعت در \en{ED}.
\item بیمار با اپیستاکسی قدامی قابل ترخیص است.
\item \en{ENT} فالوآپ ۴۸–۷۲ ساعت.
\item بیمار روی وارفارین که \en{INR} در رنج مورد انتظار است:
  ادامهٔ مصرف وارفارین.
\item \en{NSAID} حداقل ۳–۴ روز قطع شود.
\item اگر \en{pack} قرار گرفته است بیش از \en{48 h} بماند:
  \en{ABT} برای کاوریج استاف؛ آموکسی‌کلاو.
\item اگر بیمار \en{Posterior pack} شد: بستری.
\end{itemize}
\epiprint
{\footnotesize شمارهٔ دست‌نویس: ۵.\par
لبهٔ پایین تصویر، از صفحهٔ زیرین: ۴؛ باقی نوشته بریده و ناخواناست.\par}

\ifnum\EpistaxisStandalone=1
  \let\EpistaxisEnd\enddocument
\else
  \let\EpistaxisEnd\relax
\fi
\EpistaxisEnd
